\documentclass[10pt,a4paper]{amsart}
\usepackage[margin=19mm]{geometry}
\usepackage{amsmath,amssymb,mathtools}
\usepackage[scr=rsfs,cal=euler]{mathalfa}
\usepackage{xfrac}
\usepackage[unicode,hidelinks,bookmarks=false]{hyperref}
\newcommand{\ie}{i.e.\ }
\newcommand{\cf}{cf.\ }
\newcommand{\resp}{respectively\ }
\newcommand{\Subsection}{\subsection}
\providecommand{\nobreakdash}{\nobreak\hbox{-}}
\setlength{\emergencystretch}{2em}
\multlinegap0pt
\usepackage{bm}
\usepackage{bbm}
\usepackage{dsfont}
\usepackage{xspace}
\usepackage{cancel}
\usepackage{mathtools}
\usepackage{amsthm}
\makeatletter
\@ifundefined{theorem}{\newtheorem{theorem}{Theorem}}{}
\makeatother
\newcommand{\fF}{{\mathcal F}}
\title[Local $h$-polynomials, uniform triangulations and real-roote]{Local $h$-polynomials, uniform triangulations and real-rootedness}
\author{Christos A. Athanasiadis}
\date{Source: arXiv:2402.06219v2 (CC BY 4.0)}
\begin{document}
\maketitle

\noindent\emph{Source and licence.} Local $h$-polynomials, uniform triangulations and real-rootedness. Version arXiv:2402.06219v2 (CC BY 4.0), distributed under the Creative Commons Attribution 4.0 International licence, as recorded in the arXiv licence metadata for this exact version and shown on the abstract page https://arxiv.org/abs/2402.06219v2. Full rights notice: licenses/arxiv-2402.06219-CC-BY-4.0.txt.\par\smallskip

\noindent\textit{Editorial scope.} Counted unit: the Theorem whose source label is \texttt{thm:uniform} in the article (source0.tex lines 802--822, source label \texttt{thm:uniform}), delivered together with the article's own complete original proof of it (source0.tex lines 824--884, one \texttt{proof} environment of 61 lines). No mathematics is written, repaired or re-derived: statement and proof are the article's own text line for line. Required context retained uncounted: eq:h-localh (lines 636--639); eq:h-uniform (lines 719--722). Every definition, display and label of those lines is retained unchanged. Omitted from this package and not redistributed: 1--635, 640--718, 723--801, 823--823, 885--1950. Nothing inside a retained excerpt is omitted or altered: every retained body is delivered exactly as the article's lines are written, so no omission receipt of any kind accompanies this package. Citations: the retained excerpts carry no citation command at all, so no bibliography entry of the article is transcribed and no reference is redistributed. Reference adaptation: none was needed and none was made; every reference inside the delivered unit resolves inside it, so no unresolved reference or citation is produced. Printed slips kept literal: no printed word, symbol, hypothesis or number of the counted statement or of its proof was altered. Layout and class: the article's own class is \texttt{amsart} with packages including latexsym, amsmath, amssymb, epsfig, amsthm, graphicx, tikz, youngtab, hyperref; the wrapper is the lane's pinned \texttt{amsart} editorial wrapper at 10pt on A4, which restates the theorem-like environments the retained text needs and adds the article's own macro definitions that the retained text needs (\texttt{\textbackslash fF}), with extra wrapper packages (bm, bbm, dsfont, xspace, cancel, mathtools). The delivered numbering is the unit's own. Assets: no retained span contains a figure environment, a graphics-inclusion command, a PDF-page inclusion, a table or a TikZ picture; no image file is copied into this package, the package declares no asset dependency and no image file is redistributed with it. Licence: version arXiv:2402.06219v2 of this article is distributed under the Creative Commons Attribution 4.0 International licence, as recorded in the arXiv licence metadata for this exact version and shown on the abstract page https://arxiv.org/abs/2402.06219v2. A full rights notice is delivered at licenses/arxiv-2402.06219-CC-BY-4.0.txt. Characters: every retained excerpt is pure ASCII, so the delivered engine, XeLaTeX with its default Latin Modern text font, reports no missing character.
\begin{equation} \label{eq:h-localh}
  h (\Gamma, x) \, = \sum_{F \subseteq V} 
  \ell_F (\Gamma_F, x).
\end{equation}

\begin{equation} \label{eq:h-uniform}
h(\Delta,x) = \sum_{k=0}^m h_k(\Gamma)  
                        p_{\fF,m,k}(x) 
\end{equation}

\begin{theorem} \label{thm:uniform} 
Let $\Gamma$ be any triangulation of the 
$(n-1)$-dimensional simplex $2^V$. Then, there exist 
nonnegative integers $c_{k,j}(\Gamma)$ for $k, j \in 
\{0, 1,\dots,n\}$ with $k + j \le n$ such that 
%
\begin{equation} \label{eq:local-uniform}
\ell_V(\Delta,x) = \sum_{k=0}^n 
\sum_{j=0}^{n-k} c_{k,j}(\Gamma) \ell_{\fF,n,k,j}(x)
\end{equation}
%
for every $\fF$-uniform triangulation $\Delta$ of 
$\Gamma$. Specifically,
%
\begin{equation} \label{eq:ckj}
c_{k,j}(\Gamma) = \sum_{F \subseteq V: \, |F|=n-k} 
                  [x^j] \ell_F(\Gamma_F,x) 
\end{equation}
%
for $k, j \in \{0, 1,\dots,n\}$ with $k + j \le n$.
\end{theorem}

\begin{proof}
Since $\Delta_G$ is an $\fF$-uniform triangulation
of $\Gamma_G$ for every $G \subseteq V$, applying 
Equations~(\ref{eq:h-localh}) 
and~(\ref{eq:h-uniform}) we get 

\begin{align*}
\ell_V(\Delta,x) & = \sum_{G \subseteq V} 
     (-1)^{n - |G|} \, h (\Delta_G, x) \ = 
		 \sum_{G \subseteq V} (-1)^{n - |G|} 
     \sum_{j=0}^{|G|} h_j(\Gamma_G) p_{\fF,|G|,j}(x) 
		\\ & = \sum_{0 \le j \le i \le n} (-1)^{n-i} 
		       p_{\fF,i,j}(x)
    \sum_{G \subseteq V: \, |G|=i} h_j(\Gamma_G).  
\end{align*}

\medskip
\noindent
We now use Equation~(\ref{eq:h-localh}) to express 
$h_j(\Gamma_G)$ as a sum of local $h$-vector entries 
of restrictions of $\Gamma$ to faces $F \subseteq G$
of $2^V$ and change the order of summation to get 

\begin{align*}
\ell_V(\Delta,x) & = 
\sum_{0 \le j \le i \le n} (-1)^{n-i} p_{\fF,i,j}(x)
\sum_{F \subseteq G \subseteq V: \, |G|=i} [x^j] 
      \ell_F(\Gamma_F,x) \\
			& = \sum_{F \subseteq V} \sum_{j=0}^{|F|} \,
			[x^j] \ell_F(\Gamma_F,x) 
			\sum_{i=j}^n \, (-1)^{n-i} 
			{n-|F| \choose n-i} p_{\fF,i,j}(x).
\end{align*}

\medskip
\noindent
Setting $|F| = n-k$, the resulting expression can be
rewritten as 

\begin{align*}
\ell_V(\Delta,x) & = 
\sum_{0 \le j \le n-k \le n} \,
\sum_{F \subseteq V: \, |F|=n-k} [x^j] 
                                 \ell_F(\Gamma_F,x) 
			\sum_{i=n-k}^n \, (-1)^{n-i} 
			{k \choose n-i} p_{\fF,i,j}(x) \\ & = 
\sum_{0 \le j \le n-k \le n} \left( \,
\sum_{F \subseteq V: \, |F|=n-k} [x^j] 
                     \ell_F(\Gamma_F,x) \right) 
			\sum_{i=0}^k \, (-1)^i {k \choose i} 
			p_{\fF,n-i,j}(x) \\ & = 
			\sum_{0 \le j \le n-k \le n} \left( \,
\sum_{F \subseteq V: \, |F|=n-k} [x^j] 
        \ell_F(\Gamma_F,x) \right) 
				\ell_{\fF,n,k,j}(x)
\end{align*}

\medskip
\noindent
and the proof follows.
\end{proof}

\end{document}
